\documentclass[10pt,a4paper]{amsart}
\usepackage[margin=19mm]{geometry}
\usepackage{amsmath,amssymb,mathtools}
\usepackage[scr=rsfs,cal=euler]{mathalfa}
\usepackage{xfrac}
\usepackage[unicode,hidelinks,bookmarks=false]{hyperref}
\newcommand{\ie}{i.e.\ }
\newcommand{\cf}{cf.\ }
\newcommand{\resp}{respectively\ }
\newcommand{\Subsection}{\subsection}
\providecommand{\nobreakdash}{\nobreak\hbox{-}}
\setlength{\emergencystretch}{2em}
\multlinegap0pt
\usepackage{amssymb}
\usepackage{enumerate}
\newtheorem{thm}{Theorem}
\title{Pre-Schwarzian norm estimate and characterization of certain harmonic mappings}
\author{Sushil Pandit}
\date{Source: arXiv:2606.19015v1 (CC BY 4.0)}
\begin{document}
\maketitle

\noindent\emph{Source and licence.} Pre-Schwarzian norm estimate and characterization of certain harmonic mappings. Version arXiv:2606.19015v1 (CC BY 4.0), distributed by arXiv under the Creative Commons Attribution 4.0 International licence, http://creativecommons.org/licenses/by/4.0/, as recorded in the arXiv licence metadata for this exact version and shown on the abstract page https://arxiv.org/abs/2606.19015v1. Full licence notice: \texttt{licenses/arxiv-2606.19015-CC-BY-4.0.txt}.\par\smallskip

\noindent\textit{Editorial scope.} Counted unit: the article's thm thm (source0.tex lines 214--216). The article's complete original proof of it is delivered with it (source0.tex lines 217--265, one \texttt{proof} environment of 49 lines). No mathematics is written, repaired or re-derived: the statement and the proof are the article's own lines, in the article's own notation, and no retained line is substituted or re-flowed.  Retained uncounted as the source-local context the proof needs: lines 157--159.  Nothing was omitted from any retained span, so the package carries no omission record.  No retained line contains a citation, so the wrapper carries no bibliography and no bibliography entry.  No retained line contains a figure environment, an image-inclusion command or drawing code, so no image is delivered and the package declares no asset dependency.  Editorial formatting only, in addition: an AMS \texttt{amsart} preamble that restates the article's own declarations and macros used by the retained lines, namely the environments \texttt{thm} on separate counters, together with the standard packages those lines need.  The retained environments print in this document as Theorem 1 in order of appearance, whereas the article prints the same items with the numbering of its own full document.  The complete native source of the article as distributed in the arXiv e-print for this exact version is bundled unchanged as \texttt{sources/203108/source0.tex}.
    \begin{align}\label{P-151}
		h(z)=z+\sum\limits_{n=2}^\infty a_nz^n,~g(z)=\sum\limits_{n=1}^\infty b_nz^n.
	\end{align}

	\begin{thm}
		Let $f=h+\overline{g}\in\mathcal{H_R}\in\mathcal{H}$ be a harmonic mapping of the form \eqref{P-151}. Then pre-Schwarzian norm $\|P_f\|\le 3$. The estimate is sharp. 
	\end{thm}

	\begin{proof}
		The pre-Schwarzian norm $\|P_f\|$ of a harmonic mapping $f=h+\overline{g}\in\mathcal{H_R}$ is 
		\begin{align}\label{HM-100}
			\|P_f\|= & \sup_{z \in \mathbb{D}}\left|\frac{h''}{h'}-\frac{\overline{\omega}\omega'}{1-|\omega|^2}\right|(1-|z|^2)\\\nonumber
			\le & \sup_{z \in \mathbb{D}}\left|\frac{h''}{h'}\right|(1-|z|^2)+\sup_{z \in \mathbb{D}}\left|\frac{\overline{\omega}\omega'}{1-|\omega|^2}\right|(1-|z|^2)\\\nonumber
			\le & 	\|P_h\|+1.
		\end{align}
		
		
		
		
		By the hypothesis, it follows that 
		$h'(z)=\frac{m}{m-z^m},~m\in\mathbb{N}$
		and so the pre-Scwarzian derivative of $h$ is
		\begin{align*}
			P_h(z)=\frac{mz^{m-1}}{m-z^m}. 
		\end{align*}
		Now the pre-Schwarzian norm $\|P_h\|$ of $h,$ using Schwarz pick lemma, is 
		\begin{align*}
			\|P_h\|=&\sup\limits_{z\in\mathbb{D}}|P_h(z)|(1-|z|^2)\\
			=&\sup\limits_{z\in\mathbb{D}}\left|\frac{z^{m-1}}{1-z^m/m}\right|(1-|z|^2)\\
			\le &\sup\limits_{z\in\mathbb{D}}\frac{|z^{m-1}|}{1-\left|z^m/m\right|}(1-|z|^2)\\
			=&\sup\limits_{z\in\mathbb{D}}\frac{|\left(z^m/m\right)'|\left(1+\left|z^m/m\right|\right)}{1-\left|z^m/m\right|^2}(1-|z|^2)\\
			\le&\sup\limits_{z\in\mathbb{D}}\frac{1+\left|z^m/m\right|}{(1-|z|^2)}(1-|z|^2)\\
			\le& \sup\limits_{z\in\mathbb{D}}(1+|z|)\\
			=&2.
		\end{align*}
		Thus from \eqref{HM-100}, the pre-Schwarzian norm $\|P_f\|$ is
		$$\|P_f\|\le 3.$$
		
		
		To show that the estimate is sharp, we consider a harmonic mapping $f_q=h+\overline{g}$ with $h(z)=-\log{(1-z)}$ and dilatation $\omega_q(z)=\frac{z-q}{1-qz},~q\in(0,1).$ It is clear that $f_q\in \mathcal{H_R}$ and so $\|P_{f_q}\|\le3.$ A few calculations show that 
		\begin{align}
			P_h(z)=\frac{1}{1-z}\quad\text{and}\quad \frac{\overline{\omega_q(z)}\omega_q'(z)}{1-|\omega_q(z)|^2}=\frac{\overline{z}-q}{(1-qz)(1-|z|^2)}
		\end{align}
		from which it follows that 
		\begin{align}
			\|P_{f_q}\|=&\sup_{z \in \mathbb{D}}\left|\frac{h''(z)}{h'(z)}-\frac{\overline{\omega_q(z)}\omega_q'(z)}{1-|\omega_q(z)|^2}\right|(1-|z|^2)\\\nonumber
			=&\sup_{z \in \mathbb{D}}\left|\frac{1}{1-z}-\frac{\overline{z}-q}{(1-qz)(1-|z|^2)}\right|(1-|z|^2)\\\nonumber
			\ge & S_q
		\end{align}
		where 
		\begin{align}
			S_q=&\sup_{z \in [0,1)}\left|\frac{1}{1-z}-\frac{\overline{z}-q}{(1-qz)(1-|z|^2)}\right|(1-|z|^2)\\\nonumber\\\nonumber
			=&\sup_{r \in [0,1)}\left|1+r-\frac{r-q}{1-qr}\right|\\\nonumber
			=&\sup_{r \in [0,1)}\phi(r)\\\nonumber
		\end{align}
		with $\phi(r)=1+r-\frac{r-q}{1-qr}.$ We note that $\phi(0)=1+q,~\phi(1^-)=1.$ Next $\phi'(r)=0$ implies that $$r=\frac{1\pm\sqrt{1-q^2}}{q}.$$ We consider $\frac{1-\sqrt{1-q^2}}{q}=r_0$ as it lies in $[0,1).$ Since $r_0$ is stationary point and $$\phi(r_0)=\frac{2+q-2\sqrt{1-q^2}}{q}\rightarrow 3$$ as $q$ tends to $1,$ it follows that $S_q=3.$ Thus we get the relation $$3=S_q\le \|P_{f_q}\|\le 3$$ which implies that $\|P_{f_q}\|=3.$
	\end{proof}

\end{document}
